\chapter{METODE PENELITIAN}
\label{bab:metode}

\section{Alur Penelitian}

\section{Sumber Data dan Praproses}
% Korpus regulasi perpajakan; isu OCR O->0 pada nomor pasal.

\section{Ekstraksi dan Pemodelan Graf}
% Skema Neo4j: Regulation, Article/Pasal, Concept; relasi BELONGS_TO,
% REFERENCES, AMENDS, DEFINES. Format ID pasal "<regulation_id>::<article_number>".

\section{Tahap Seeding Bersama}
% Kunci komparabilitas: kedua sistem melalui src/seeding.py yang sama
% (dense ChromaDB pool + BM25 pool, difusikan dengan RRF). Setiap perbedaan
% terukur berasal dari tahap graf saja.

\section{Sistem 1 --- Baseline RAG (Hibrida)}

\section{Sistem 2 --- GraphRAG (Ekspansi Graf)}

\section{Rancangan Evaluasi}
% Ground-truth eval.jsonl; pembagian train/held-out dibekukan SEBELUM tuning
% TOP_K_VECTOR / GRAPH_HOP_DEPTH. Metrik: recall@k, dsb. Harness 2x2.

\section{Lingkungan Implementasi}
% Python, ChromaDB, Neo4j, Ollama (qwen3.5:9b, qwen3-embedding:0.6b).
